\chapter{Top-Down Parsing}

\section{Idea}

A \textbf{top-down} parser builds the parse tree from the root down, producing a
\emph{leftmost derivation}. At every step it must decide which production to apply to the
leftmost non-terminal. An LL(1) parser makes that decision by peeking at \emph{one} token of
input. (``LL(1)'': read input \textbf{L}eft to right, produce a \textbf{L}eftmost derivation,
\textbf{1} token of lookahead.)

We augment every grammar with a new start production $S' \to\ \vdash S \dashv$, where
$\vdash$ (BOF) and $\dashv$ (EOF) are new terminals surrounding the input. Then there is
always a lookahead symbol, even at the end.

\section{Stack-based top-down parsing}

The parser keeps a stack of grammar symbols representing the part of the derivation not yet
matched. Invariant: \emph{consumed input} $+$ \emph{stack (read top to bottom)} $=$ current
sentential form of the leftmost derivation.

\begin{lstlisting}[style=cpp]
push S';                        // or push the right-hand side of S' in reverse
for each token a in input (including BOF, EOF):
    while top of stack is a non-terminal A:
        if Predict[A][a] is empty: reject;
        pop A; push the RHS of Predict[A][a] in reverse order;
    if top of stack != a: reject;
    pop;                        // match
accept if stack is empty and input is exhausted;
\end{lstlisting}

Everything depends on the table $\mathrm{Predict}[A][a]$: the production to use when $A$ is on
top of the stack and $a$ is the next input token. A grammar is \textbf{LL(1)} iff every entry
contains at most one production.

\section{Nullable, FIRST and FOLLOW}

\begin{definition}
For $\beta\in V^*$ and $A\in N$:
\begin{itemize}
  \item $\mathrm{Nullable}(\beta)$ is true iff $\beta \Rightarrow^* \varepsilon$.
  \item $\mathrm{First}(\beta) = \{ a\in T : \beta \Rightarrow^* a\gamma \}$: terminals that can
        begin a string derived from $\beta$.
  \item $\mathrm{Follow}(A) = \{ b \in T : S' \Rightarrow^* \alpha A b \gamma\}$: terminals that
        can appear immediately after $A$ in some sentential form.
\end{itemize}
\end{definition}

\begin{definition}
$\mathrm{Predict}(A, a) = \{A\to\beta : a\in\mathrm{First}(\beta)\} \cup
\{A\to\beta : \mathrm{Nullable}(\beta) \text{ and } a \in \mathrm{Follow}(A)\}$.
\end{definition}

\subsection{Computing them (fixed-point iteration)}

All three are computed by iterating until nothing changes.

\begin{lstlisting}[style=cpp]
// Nullable
initialize Nullable[A] = false for all A
repeat
  for each A -> B1 ... Bk:
    if k == 0 or all Bi are nullable non-terminals: Nullable[A] = true
until no change

// First (for non-terminals)
initialize First[A] = {} for all A
repeat
  for each A -> B1 ... Bk:
    for i = 1..k:
      if Bi is a terminal: First[A] += {Bi}; break
      First[A] += First[Bi]
      if not Nullable[Bi]: break
until no change

// Follow
initialize Follow[A] = {} for all A
repeat
  for each A -> B1 ... Bk:
    for i = 1..k with Bi a non-terminal:
      Follow[Bi] += First(B(i+1) ... Bk)
      if Nullable(B(i+1) ... Bk): Follow[Bi] += Follow[A]
until no change
\end{lstlisting}

For a sequence: $\mathrm{First}(X_1\dots X_k) = \mathrm{First}(X_1) \cup \mathrm{First}(X_2)$
(if $X_1$ nullable) $\cup\ \dots$, stopping at the first non-nullable symbol; for a terminal
$a$, $\mathrm{First}(a) = \{a\}$.

\section{A worked example}

The unambiguous expression grammar from Chapter 8 is left-recursive and therefore not LL(1)
(see below). Here is an equivalent LL(1) version:
\begin{center}
\begin{tabular}{rl}
0 & $S' \to\ \vdash E \dashv$ \\
1 & $E \to T E'$ \\
2 & $E' \to + T E'$ \\
3 & $E' \to \varepsilon$ \\
4 & $T \to F T'$ \\
5 & $T' \to * F T'$ \\
6 & $T' \to \varepsilon$ \\
7 & $F \to ( E )$ \\
8 & $F \to \mathtt{id}$
\end{tabular}
\end{center}

\begin{center}
\begin{tabular}{llll}
\toprule
 & Nullable & First & Follow \\
\midrule
$E$  & no  & $\{(, \mathtt{id}\}$ & $\{), \dashv\}$ \\
$E'$ & yes & $\{+\}$ & $\{), \dashv\}$ \\
$T$  & no  & $\{(, \mathtt{id}\}$ & $\{+, ), \dashv\}$ \\
$T'$ & yes & $\{*\}$ & $\{+, ), \dashv\}$ \\
$F$  & no  & $\{(, \mathtt{id}\}$ & $\{*, +, ), \dashv\}$ \\
\bottomrule
\end{tabular}
\end{center}

Predict table (entries are rule numbers; blank = error):
\begin{center}
\begin{tabular}{c|ccccccc}
 & $\vdash$ & $\dashv$ & id & $+$ & $*$ & $($ & $)$ \\
\hline
$S'$ & 0 & & & & & & \\
$E$  & & & 1 & & & 1 & \\
$E'$ & & 3 & & 2 & & & 3 \\
$T$  & & & 4 & & & 4 & \\
$T'$ & & 6 & & 6 & 5 & & 6 \\
$F$  & & & 8 & & & 7 & \\
\end{tabular}
\end{center}
No entry has two rules, so the grammar is LL(1).

\begin{example}[Parsing $\vdash \mathtt{id} + \mathtt{id} * \mathtt{id} \dashv$]
Stack shown with top on the left.
\begin{center}
\small
\begin{tabular}{rlll}
\toprule
step & stack & remaining input & action \\
\midrule
1 & $S'$ & $\vdash$ id + id * id $\dashv$ & expand 0 \\
2 & $\vdash E \dashv$ & $\vdash$ id + id * id $\dashv$ & match $\vdash$ \\
3 & $E \dashv$ & id + id * id $\dashv$ & expand 1 \\
4 & $T E' \dashv$ & id + id * id $\dashv$ & expand 4 \\
5 & $F T' E' \dashv$ & id + id * id $\dashv$ & expand 8 \\
6 & id $T' E' \dashv$ & id + id * id $\dashv$ & match \\
7 & $T' E' \dashv$ & + id * id $\dashv$ & expand 6 \\
8 & $E' \dashv$ & + id * id $\dashv$ & expand 2 \\
9 & $+ T E' \dashv$ & + id * id $\dashv$ & match \\
10 & $T E' \dashv$ & id * id $\dashv$ & expand 4 \\
11 & $F T' E' \dashv$ & id * id $\dashv$ & expand 8 \\
12 & id $T' E' \dashv$ & id * id $\dashv$ & match \\
13 & $T' E' \dashv$ & * id $\dashv$ & expand 5 \\
14 & $* F T' E' \dashv$ & * id $\dashv$ & match \\
15 & $F T' E' \dashv$ & id $\dashv$ & expand 8 \\
16 & id $T' E' \dashv$ & id $\dashv$ & match \\
17 & $T' E' \dashv$ & $\dashv$ & expand 6 \\
18 & $E' \dashv$ & $\dashv$ & expand 3 \\
19 & $\dashv$ & $\dashv$ & match \\
20 & (empty) & (empty) & accept \\
\bottomrule
\end{tabular}
\end{center}
The sequence of expansions (0,1,4,8,6,2,4,8,5,8,6,3) is exactly the leftmost derivation.
\end{example}

The parse tree the trace builds, root first (leaves read left to right give the input):
\begin{center}
\begin{tikzpicture}[level distance=0.8cm,every node/.style={font=\small},
  level 1/.style={sibling distance=3.5cm},
  level 2/.style={sibling distance=3cm},
  level 3/.style={sibling distance=1.5cm},
  level 4/.style={sibling distance=1.6cm},
  level 5/.style={sibling distance=0.9cm}]
\node {$S'$}
  child {node {$\vdash$}}
  child {node {$E$}
    child {node {$T$}
      child {node {$F$} child {node {id}}}
      child {node {$T'$} child {node {$\varepsilon$}}}}
    child {node {$E'$}
      child {node {$+$}}
      child {node {$T$}
        child {node {$F$} child {node {id}}}
        child {node {$T'$}
          child {node {$*$}}
          child {node {$F$} child {node {id}}}
          child {node {$T'$} child {node {$\varepsilon$}}}}}
      child {node {$E'$} child {node {$\varepsilon$}}}}}
  child {node {$\dashv$}};
\end{tikzpicture}
\end{center}

\section{When a grammar is not LL(1)}

\subsection{Left recursion}
If $A \to A\alpha \mid \beta$, then $\mathrm{First}(\beta) \subseteq \mathrm{First}(A\alpha)$,
so both rules are predicted on the same token: a conflict. (A recursive-descent parser would
loop forever.) \textbf{Fix}: rewrite as right recursion with a new non-terminal:
\[
A \to A\alpha \mid \beta \qquad\leadsto\qquad A \to \beta A', \quad A' \to \alpha A' \mid \varepsilon.
\]
This is how $E \to E + T \mid T$ became $E\to TE'$, $E' \to +TE' \mid \varepsilon$.

\begin{important}[Associativity is lost]
The rewritten grammar's parse tree for $a-b-c$ hangs to the right. A compiler must
re-associate when building the abstract syntax tree (or evaluate with an accumulator) to
preserve left associativity. This is one reason LR parsing is preferred for expressions.
\end{important}

\begin{center}
\begin{tikzpicture}[level distance=0.8cm,level 1/.style={sibling distance=2.2cm},
  level 2/.style={sibling distance=1.4cm},level 3/.style={sibling distance=0.85cm},every node/.style={font=\small}]
\node {$E$}
  child {node {$T$} child {node {$a$}}}
  child {node {$E'$}
    child {node {$-$}}
    child {node {$T$} child {node {$b$}}}
    child {node {$E'$}
      child {node {$-$}}
      child {node {$T$} child {node {$c$}}}
      child {node {$E'$} child {node {$\varepsilon$}}}}};
\end{tikzpicture}

\small $a-b-c$ with $E\to TE'$, $E'\to -TE'\mid\varepsilon$: the tree hangs to the right,
so a naive walk would compute $a-(b-c)$.
\end{center}

\subsection{Common prefixes: left factoring}
$A \to \alpha\beta_1 \mid \alpha\beta_2$ conflicts on $\mathrm{First}(\alpha)$. Factor out
the prefix: $A \to \alpha A'$, $A' \to \beta_1 \mid \beta_2$. E.g.\ WLP4's
\code{factor -> ID | ID LPAREN RPAREN | ...} needs factoring for LL(1).

\subsection{Limits}
Some languages have no LL(1) grammar at all (e.g.\ $\{a^nb^m : n\ge m\}$ which needs
unbounded lookahead to decide). LL($k$) for larger $k$ helps sometimes but tables grow
quickly. Every LL(1) grammar is unambiguous.

\section{Recursive descent}

An LL(1) grammar can be parsed by hand by writing one function per non-terminal; the call
stack plays the role of the parse stack.
\begin{lstlisting}[style=cpp]
Token la;                                  // lookahead
void match(Kind k) { if (la.kind != k) error(); la = next(); }
void E()  { T(); Eprime(); }
void Eprime() {
    if (la.kind == PLUS) { match(PLUS); T(); Eprime(); }
    else if (la.kind == RPAREN || la.kind == EOF_) { /* epsilon */ }
    else error();
}
void T()  { F(); Tprime(); }
void Tprime() {
    if (la.kind == STAR) { match(STAR); F(); Tprime(); }
    else if (la.kind == PLUS || la.kind == RPAREN || la.kind == EOF_) {}
    else error();
}
void F() {
    if (la.kind == ID) match(ID);
    else if (la.kind == LPAREN) { match(LPAREN); E(); match(RPAREN); }
    else error();
}
\end{lstlisting}
Each branch condition is exactly a row of the predict table. Many production compilers
(GCC's and Clang's C++ front ends) use hand-written recursive descent.

\section{Summary}
\begin{itemize}
  \item Top-down parsing = leftmost derivation; expand non-terminals on the stack, match terminals.
  \item Predict$(A,a)$ uses First of the RHS, plus Follow$(A)$ when the RHS is nullable.
  \item Compute Nullable, First, Follow by fixed-point iteration.
  \item LL(1) iff the predict table has no multiply-defined entries.
  \item Remove left recursion and left-factor to make a grammar LL(1); beware lost associativity.
  \item Recursive descent is a hand-coded LL(1) parser.
\end{itemize}
